\section{Bridging Blockchains}

As the DeFi space continues to experience exponential growth, the number of blockchains available for investment also increases. However, each blockchain has its characteristics, with some being easy to use but expensive in fees, while others have light fees but are complex to navigate. The Zoo Bridge Protocol addresses these limitations by developing a bridge between major blockchains, facilitated by \$ZOO. This bridge, established through the Zoo platform, facilitates universal exchange within the blockchain space.

\subsection{Threshold Signature Scheme}

While traditional bridges trap tokens in vault contracts, the Zoo Bridge Protocol implements a unique approach, leveraging a threshold signature scheme. Oracles monitor the blockchain network to ensure funds are deposited in the smart contract, with ZOO tokens being burned on one chain via the bridge protocol and minted on the desired chain. This eliminates the need for a vault contract that could potentially contain vulnerabilities that could be exploited by attackers.

The Zoo Bridge is entirely secured by cross-chain multi-party computation (MPC) validators who operate with threshold signature schemes (TSS).

\subsection{Leaderless Consensus}

The network is leaderless and ensures security through each validator running the same process upon receiving on-chain events on the various chains that the MPC validator group tracks. Consensus is achieved when two-thirds of all validators collectively sign the same transaction using their key, and a trade is issued to the destination chain. The \$ZOO is minted when enough oracles see the deposit into the smart contract. In a standard transaction, such as when ZOO teleports from one chain to another, three network confirmations are awaited.

The Zoo Bridge Protocol represents a significant innovation in blockchain technology, improving the interoperability of blockchains and expanding investment opportunities.
